\ifx\danslelivre\undefined
\documentclass[10pt]{../fichiers-configuration-latex/smfart}
\usepackage{ifthen}
\newboolean{english}
\newboolean{francais}
\setboolean{english}{false}
\setboolean{francais}{true}
\input{../fichiers-configuration-latex/paquets}
%\input{../fichiers-configuration-latex/paquets-isolatin}
\selectlanguage{francais}
\input{../fichiers-configuration-latex/lemme-proposition-theoreme}
%\input{../fichiers-configuration-latex/lemme-proposition-theoreme-ascii}
\input{../fichiers-configuration-latex/numerotation}
\input{../fichiers-configuration-latex/macros}
%\input{../.cv}

\externaldocument{../01-quasi-excellents/excellent}
\externaldocument{../02-topologies/topologies}

\begin{document}


\author{Luc Illusie et Yves Laszlo}


\title{III. Approximation}
%\date{exposés oraux des  11 et 18 juin 2008, notes du \today}

\setcounter{tocdepth}{1}
\maketitle
%\version
\tableofcontents

\else
\setboolean{english}{false}
\setboolean{francais}{true}
\selectlanguage{francais}

\chapter{Approximation}
\addtocontents{toc}{par Luc Illusie et Yves Laszlo}
\begin{center}
Luc Illusie et Yves Laszlo
\end{center}
\fi

\begin{abstract}On montre ici comment des méthodes
d'approximation de complexes et le théorème de Popescu
(\refext{I}{1-Artin-Popescu}) permettent de réduire la preuve
du théorème d'uniformisation des schémas quasi-excellents
(\ref{3-unif}) au cas complet.\end{abstract}



\section{Introduction} On montre ici comment ramener la preuve
du théorème d'uniformisation (\ref{3-unif}) au cas local,
noethérien complet (\ref{3-reduc}). On utilise pour cela
le théorème de Popescu (qui implique que les anneaux locaux
noethériens, henséliens et excellents vérifient la propriété
d'approximation d'Artin, cf.~\refext{I}{1-Artin-Popescu}) et des
méthodes d'approximations de complexes de longueur $2$ adaptées
de \cite{Con-deJ02} (cf. section~\ref{3-section-appro}).

L'exposé oral donné par Alban Moreau utilisait des résultats
(dus à Ofer Gabber) d'approximations de complexes plus forts
que ceux utilisés ici (\ref{3-CdJ}). Une version écrite de son
exposé a été très utile pour la rédaction de ce texte : nous
l'en remercions. Nous remercions également Fabrice Orgogozo de
nous avoir signalé que l'énoncé \cite[3.1]{Con-deJ02} suffisait
pour les applications en vue.

\section{Modèles et approximations à la Artin-Popescu}\label{3-mod}




Soit $A$  un anneau local noethérien, $\m$ son idéal maximal,
$\hat A$ son complété. On suppose $A$ excellent et hensélien. Soit
$\pi:{\s}=\Spec(\hat A)\ra S=\Spec(A)$ le morphisme canonique.  Pour
tout $n\geq 0$, on  note $$i_n:S_n\hookrightarrow {\s}$$ l'immersion
fermée définie par d'idéal $\hat \m^{n+1}=\m^{n+1}\hat A$ de
$\hat A$. Le composé $$\pi i_n:S_n\ra{\s}\ra S$$ est l'immersion
fermée $S_n\hookrightarrow S$ définie par l'idéal $\m^{n+1}$.

\begin{definition}  Soient $g:{\s}\ra T$ et $f:S\ra T$ des
morphismes de schémas et $n\in\NN$. On dira que  $f$ et $g$
sont $(n+1)$-proches si leurs restrictions $f\pi i_n$ et $gi_n$
à $S_n$ coïncident.  \end{definition}



Si $\X$ est un ${\s}$-schéma, on note $\X_n$ le $S_n$-schéma
$\X\times_{{\s}}S_n\ra S_n$.


Écrivons $\hat A$ comme limite inductive suivant
un ensemble ordonné filtrant $E$ de $A$-algèbres
de type fini $A_\alpha$. On a  des diagrammes commutatifs
\begin{equation}\label{3-salpha}\xymatrix{&S_\alpha=\Spec(A_\alpha)\ar[d]^{t_\alpha}\\{\s}\ar[ru]^{s_\alpha}\ar[r]^\pi&S}\end{equation}
avec  $t_\alpha$ de type fini et un isomorphisme  ${\s}=\varprojlim
S_\alpha$ \ega{IV$_3$}{8.2.3}.


\begin{definition}\label{3-def-mod} Soit $\X$ un
${\s}$-schéma de type fini et $h:\X\ra\Y$ un morphisme
de $\s$-schémas de type fini.  \begin{itemize}\item[]
Un modèle de $\X$  sur $S_\alpha$ est un diagramme cartésien
$$\xymatrix{\X\ar[r]\ar[d]_f\ar@{}[rd]|\square&X_\alpha\ar[d]^{f_\alpha}\\{\s}\ar[r]&S_\alpha\\}$$
où $X_\alpha$ est de type fini sur $S_\alpha$.  \item[] Un
modèle de $h$ sur $X_\alpha$ est un $S_\alpha$-morphisme
$h_\alpha:X_\alpha\ra Y_\alpha$ muni d'un isomorphisme
$h\iso(h_\alpha)_{{\s}}$.\end{itemize} \end{definition}

Des modèles de $\X$ sur $S_\alpha$  existent pourvu que
$\alpha$ soit assez grand \ega{IV$_3$}{8.8.3}. De plus, si
$X_\alpha,X_\beta$ sont des modèles de $\X$ sur $S_\alpha,S_\beta$,
il existe $\gamma\geq \alpha,\beta$ et un $S_{\gamma}$-isomorphisme
$$X_\alpha\times_{S_\alpha}S_\gamma\xrightarrow{\sim}
X_\beta\times_{S_\beta}S_\gamma$$ (\textit{loc. cit.}). De même,
des modèles $h_\alpha$ de $h:\X\ra\Y$ sur $S_\alpha$  existent
pourvu que $\alpha$ soit assez grand et les images inverses de tels
modèles $h_\alpha,h_\beta$ sur $S_\gamma$ sont $S_\gamma$-isomorphes
pour $\gamma\geq\alpha,\beta$ assez grand.

Si $T$ est un $S$-schéma et $B$ une $A$-algèbre, on note
$T(B)=\Hom_S(\Spec(B),T)$ l'ensemble des $S$-points de $T$ à valeurs
dans $B$. D'après le théorème de Popescu \cite[1.3]{Pop86},
comme $A$ est excellent et hensélien, il vérifie la propriété
d'approximation d'Artin, cf.  \refext{I}{1-Artin-Popescu}. Donc,
comme  $S_\alpha\ra S$  est de type fini,  $S_\alpha(A)$ est dense
dans $S_\alpha(\hat A)$ (pour la topologie $\m$-adique). Il existe
donc, pour tout $n\geq 0$ une section $u:S\ra S_\alpha$ de $t_\alpha$
qui est $n$-proche de $s_\alpha:\s\ra S_\alpha$. On définit alors
$X_u$ par le diagramme cartésien

\begin{equation}\label{3-Xu}
\xymatrix{X_u\ar[r]\ar[d]_{f_u}\ar@{}[rd]|\square&X_\alpha\ar[d]^{f_\alpha}\\S\ar[r]^u&S_\alpha\\}
\end{equation}



Comme $u$ est $n$-proche de $s_\alpha$, on a par définition
l'égalité $$u\pi i_n=s_\alpha i_n$$ de sorte la restriction de
$X_u\ra S$ à $S_n$ s'identifie à $\X_n\ra S_n$, autrement dit on
a un carré cartésien \begin{equation}\label{3-Xn}
    \xymatrix{\X_n\ar[r]\ar[d]_{f_n}\ar@{}[rd]|\square&X_u\ar[d]^{f_u}\\S_n\ar[r]^{i_n}&S\\}
\end{equation}

De même, si $\X,{{\Y}}$ sont de type fini sur ${\s}$ et $h_\alpha$
est un modèle de $h\in\Hom_{{\s}}(\X,{{\Y}})$ sur $S_\alpha$,
l'image inverse $h_u:X_u\ra Y_u$ est un $S$-morphisme induisant la
restriction $h_n:\X_n\ra {{\Y}}_n$ de $h$ au dessus de $S_n$.



\section{Approximations et topologie des altérations} Commençons
par un rappel (cf. \refext{II}{2-topologies}) sur la topologie
des altérations. Soit $T$ un schéma noethérien.  La catégorie
${\pspf}/T$ est la sous-catégorie pleine de la catégorie des
$T$-schémas dont les objets sont les $T$-schémas réduits de
type fini $X$, dont tout point maximal s'envoie sur un point
maximal de $T$ avec extension résiduelle finie. Notons que les
morphismes de $\pspf/T$ envoient point maximal sur point maximal.
  On définit deux topologies sur $\pspf/T$.



\begin{itemize}
  \item[(i)] La \textit{topologie des altérations} est la
  moins fine pour laquelle les familles suivantes sont couvrantes
  \begin{itemize} \item[(a)]les recouvrements ouverts de Zariski ;
  \item[(b)] les morphismes propres et surjectifs .  \end{itemize}
Une famille couvrante pour la topologie des altérations sera dite
\pspf-couvrante.
  \item[(ii)] Soit $\ell$ un nombre premier. La \textit{topologie
  des $\ell'$-altérations} sur ${\pspf}/T$ est la moins fine pour
  laquelle les familles suivantes sont couvrantes

  \begin{itemize} \item[(a)]les recouvrements étales de Nisnevich ;
  \item[(b)] les morphismes  propres surjectifs $\X'\ra \X$ tels que
  pour tout point maximal $\eta$ de $X$, il existe un point maximal
  $\eta'$ de $X'$ au dessus de $\eta$ avec  $\ell$ ne divisant pas
  $\deg(k(\eta')/k(\eta))$.  \end{itemize} Une famille couvrante
  pour la topologie des $\ell'$-altérations
sera dite \altl-couvrante.  \end{itemize}




Pour tout $T$-schéma $X$ dominant, on note $X_\r$ le sous-schéma
fermé réduit de $X$  réunion des composantes irréductibles qui
dominent une composante irréductible de $T$.




\begin{proposition}\label{3-alt-mod}On reprend les notations de
\ref{3-mod} : soit $A$  un anneau local noethérien, $\m$ son
idéal maximal, $\hat A$ son complété. On suppose $A$ excellent
et hensélien. Soit $\pi:{\s}=\Spec(\hat A)\ra S=\Spec(A)$
le morphisme canonique.Soit $\X\ra {\s}$  un objet non vide de
${\pspf/{\s}}$. Supposons de plus $S$ intègre.
 \begin{itemize} \item[(i)]Alors, il existe $\alpha_0\in E,
 n_0\in\mathbf{N}$ tel que  pour tout $\alpha\geq\alpha_0$, tout
 entier $n\geq n_0$, toute section $u:S\ra S_\alpha$ de $t_\alpha$
 qui est $n$-proche de $s_\alpha:{\s}\ra S_\alpha$, tout modèle
 (cf. \ref{3-def-mod}) $X_\alpha$ de $\X$ sur $S_\alpha$, $X_u$ est
 à fibre générique finie et le morphisme composé $(X_u)_\r\ra
 X_u\ra S$ est un objet non vide de ${\pspf/S}$.
   \item[(ii)] Supposons $X \ra{\s}$ \pspf-couvrant. Alors,
   il existe $\alpha_0\in E, n_0\in\mathbf{N}$ tel que  pour
   tout $\alpha\geq\alpha_0$, tout entier $n\geq n_0$, toute
   section $u:S\ra S_\alpha$ de $t_\alpha$ qui est $n$-proche de
   $s_\alpha:{\s}\ra S_\alpha$, tout modèle $X_\alpha$ de $\X$
   sur $S_\alpha$, le morphisme composé $(X_u)_\r\ra X_u\ra
   S$ est $\pspf$-couvrant.  \item[(iii)] Supposons $X \ra{\s}$
   \altl-couvrant. Alors, il existe $\alpha_0\in E, n_0\in\mathbf{N}$
   tel que  pour tout $\alpha\geq\alpha_0$, tout entier $n\geq n_0$,
   toute section $u:S\ra S_\alpha$ de $t_\alpha$ qui est $n$-proche
   de $s_\alpha:{\s}\ra S_\alpha$, tout modèle $X_\alpha$ de $\X$
   sur $S_\alpha$, le morphisme composé $(X_u)_\r\ra X_u\ra S$
   est $\altl$-couvrant.
 \end{itemize}
\end{proposition}






\begin{proof} Observons d'abord que, $S$ étant
hensélien et excellent, $\s$ est intègre,
cf. \refext{I}{1-cor-normal-versus-complet}.

Prouvons (i). Comme $\X\ra\s$ est génériquement fini, il existe
$a\in\hat A-\{0\}$ tel que $\X$ soit fini, surjectif et libre
de rang $d>0$ au dessus de l'ouvert non vide  $\X-V(a)$. On peut
choisir $\alpha_0$ assez grand de sorte que \begin{itemize}
  \item $a$ provienne de $a_\alpha\in A_\alpha-\{0\}$ pour
  $\alpha\geq\alpha_0$ ; \item  $X_\alpha\ra S_\alpha$ soit fini,
  surjectif (\ega{IV$_3$}{8.10.5}) et libre de rang $d$ sur
  $X_\alpha-V(a_\alpha)$ (utiliser \ega{IV$_3$}{8.5.2}).
\end{itemize}

Choisissons alors un entier $n$ tel que $a\not\in\hat\m^{n+1}$. Pour
tout $\alpha\geq\alpha_0, m\geq n$, toute  section $u$ qui est
$m$-proche de $t_\alpha,$, on a $$u^*(a_\alpha)\not\in\m^{n+1}$$
et donc $u^*(a_\alpha)$ est non nul. Ceci assure que  $X_u$ est
fini, surjectif et libre de rang $d$ au dessus de l'ouvert non vide
$S-V(u^*(a_\alpha))$ image réciproque de de $S_\alpha-V(a_\alpha)$
par $u$. Le premier point en découle.


Prouvons (iii) [La preuve de (ii) est en tout point similaire]. On
suppose donc que  $\X\ra {\s}$ est $\altl$-couvrant. On sait
(\refext{II}{2-forme standard l'-altérations}) que $\X\ra {\s}$
est dominé dans $\pspf/\s$ par un recouvrement standard $${{\Y}}\ra
\X'\ra {\s}$$ avec

 \begin{itemize}
   \item ${{\Y}}\ra \X'$  Nisnevich couvrant \item $\X'\ra {\s}$
   propre et surjectif dont la restriction à chaque composante
   irréductible est dominante et génériquement, le degré
   générique de l'une d'elles étant premier à $\ell$.
 \end{itemize}

 Quitte à remplacer le schéma réduit $\X'$ par une  composante
 convenable et $\Y$ par le \altl-recouvrement Nisnevich induit,
 on peut supposer $\X'$ \emph{intègre  } de degré générique
 $\deg(\X'/\s)=\delta$ premier à $\ell$.

Soit $\eta$ le point générique de $S$. La construction
$X\mapsto X_\r$ est fonctorielle pour la sous-catégorie pleine des
$S$-schémas $X$ à fibre générique  finie. Or, d'après (i), pour
des choix de modèles et de section $u$ de $t_\alpha$ convenables,
on sait que $Y_u,X'_u$ et $X_u$ sont à fibre générique finie. On
a donc une factorisation
 $$\xymatrix{(Y_u)_\r\ar[r]\ar[d]&(X_u)_\r\ar[d]\\(X'_u)_\r\ar[r]&S}$$
 Or, toujours d'après (i), on peut en outre supposer que
 $(Y_u)_\r,(X'_u)_\r$ et $(X_u)_r$ sont des objets de $\pspf/S$.
 Pour conclure que $(X_u)_\r\ra S$ est $\altl$-couvrant, il
 suffit de prouver que pour $u$ convenable $(Y_u)_\r\ra S$ est
 $\altl$-couvrant.

 Tenant compte des propriétés de permanence usuelles des modèles
 \ega{IV$_3$}{8.8.3 et 8.10.5}, la preuve de (i) assure que pour
 des modèles et $u$ convenables le $X'_u\rightarrow S$ est propre
 et surjectif et que sa fibre générique est de degré premier
 à $\ell$. Ceci assure que la restriction de $X'_u\rightarrow S$
 à au moins une des composantes réduites de $X'_u$ dominant $S$
 est de degré premier à $\ell$. 
 Ainsi,  $(X'_u)_\r\ra S$ est bien \altl-couvrant.

 La propriété d'être un recouvrement Nisnevich (resp. propre et
 surjectif) étant stable par changement de base, reste à prouver
 le lemme suivant.

\begin{lemme2}\label{3-stabmod}Il existe $\alpha_0$ tel que
pour tout $\alpha\geq\alpha_0$, tout modèle ${ Y}_\alpha\ra {
X}'_\alpha\ra  S_\alpha$ de ${{\Y}}\ra \X'\ra {\s}$ vérifie
$\Y_\alpha\ra \X'_\alpha$ est Nisnevich couvrant.  \end{lemme2}

\begin{proof}Dire que ${{\Y}}\ra \X'$ est Nisnevich couvrant, c'est
dire qu'il est lisse, quasi-fini et qu'on a une stratification
$$\varnothing=\X'_0\subset \X'_1\cdots \subset \X'_n=\X'$$ avec
$\X'_i$ fermé de $\X'$ et ${{\Y}}/\X'$ a une section au dessus
de $\X'_{i+1}-\X'_i$. La conclusion découle immédiatement de
cette remarque et des propriétés de permanence usuelles des
modèles \ega{IV$_3$}{8.8.3 et 8.10.5} et \ega{IV$_4$}{17.7.8}.
\end{proof}\end{proof}



Le but de ce qui suit est d'améliorer les résultats topologiques de
la proposition~\ref{3-alt-mod} en montrant que des épaississements
convenables des cônes normaux des fibres spéciales de $\X$
(resp. $X_u$) dans $\X$ (resp. $X_u$) sont isomorphes. Ceci
permettra de prouver des énoncés de stabilité de propriétés
dans le passage de $\X$ à $X_u$, en l'occurrence la dimension et
la régularité (corollaire~\ref{3-perm}).


\section{Gradués supérieurs et approximations de
complexes}\label{3-section-appro}  Soient $I$ un idéal d'un
topos annelé $(\mathscr{X},\OO)$, $\F$ un ${\OO}$-module
de $\mathscr{X}$ et $a$ un entier $\geq 1$. On pose
$I^n=\OO$ si $n\leq 0$. On définit le module $\ZZ$-gradué
$$\gr_a(\F)=\mathop{\oplus}_{n\in\ZZ}I^n\F/I^{n+a}\F$$ qui est
donc la somme $$\gr_a(\F)=\F/{I}\F\oplus\cdots\oplus\F/{I}^aF\oplus
{I}\F/{I}^{a+1}\F\oplus {I}^2\F/{I}^{a+2}\F\oplus\cdots$$ concentrée
en degrés $\geq -(a-1)$.  C'est un ${\OO}/I^a$-module ; de plus,
le produit $$I^n\otimes I^m\ra I^{n+m}$$ induit une structure de
${\OO}/I^a$-algèbre $\ZZ$-graduée sur $\gr_a({\OO})$ et $\gr_a(\F)$
est un $\gr_a(\OO)$-module $\ZZ$-gradué.

On s'intéresse ici au cas où $\mathscr{X}$ est le topos de Zariski
d'un $S$-schéma $X$ annelé par son faisceau structural $\OO$
et $I=\hat m{\OO}$.

\begin{remarque}Le morphisme surjectif tautologique
$\gr_a(\OO)\ra\gr_1(\OO)$ a pour noyau $J=I\gr_a(\OO)$. On a donc
$J^a=0$ (puisque $J$ est un $\OO_{X_{a-1}}$-module) de sorte que
$C_a(\X)=\Spec(\gr_a(\OO))$ est un épaississement d'ordre $a-1$
du cône normal $\Spec(\gr_1(\OO))$.\end{remarque}


\begin{definition} Soient $\X,\Y$ des $S$-schémas (resp. des
$\hat S$-schémas). Un $a$-isomorphisme $\X\iso_a {{\Y}}$
est la donnée d'un $S$-isomorphisme $\phi:X_{a-1}\iso
Y_{a-1}$ et d'un isomorphisme de $\gr_a(A)$-algèbres graduées
$\phi^{-1}\gr_a(\OO_{{\Y}})\iso\gr_a(\OO_\X)$. On dit dans ce cas
que $\X,{{\Y}}$ sont $a$-proches. \end{definition}

On identifiera alors leurs fibres spéciales $X_0,Y_0$ grâce à
l'isomorphisme $X_{a-1}\iso Y_{a-1}$.



\subsection{} On adapte ici le théorème 3.2 de \cite{Con-deJ02}
(et le lemme clef 3.1 de \textit{loc. cit.}). Commençons par
une définition. Soient $B$ un anneau noethérien et $I$ un idéal
de $B$.

\begin{definition}\label{3-AR} Soit $f:M\ra N$ un morphisme de
$B$-modules de type fini. Un entier $c\geq 0$ est une constante
d'Artin-Rees de $f$ si pour tout $n\geq c$ on a $${{I}}^nN\cap
\Im(f)\subset {{I}}^{n-c}\Im(f).$$\end{definition}

Le lemme d'Artin-Rees assure l'existence d'une constante
d'Artin-Rees.


\begin{proposition}\label{3-CdJ} Soient
$(L^\bullet,\d_L^\bullet),(M^\bullet,\d_M^\bullet)$ des
complexes de $B$-modules libres de type fini concentrés
en degré $[-2,0]$ avec $L^i=M^i$ pour tout $i$. Soit $c$
une constante d'Artin-Rees pour $\d_L^{-2}$ et $\d_L^{-1}$
et $n$ un entier $\geq c$. Supposons $H^{-1}(L^\bullet)=0$ et
$$\d_L^\bullet=\d_M^\bullet\mod I^{n+1}.$$ Alors : \begin{itemize}
\item[(i)] $c$ est une constante d'Artin-Rees pour $\d_M^{-1}$
; \item[(ii)] $H^{-1}(M^\bullet)=0$ ; \item[(iii)] L'identité
de $L^0=M^0$induit un isomorphisme de $\gr_{n+1-c}(B)$-modules
$$\gr_{n+1-c}(H^{0}(L^\bullet))\iso\gr_{n+1-c}(H^{0}(M^\bullet)) ;$$
\item[(iv)] De plus, si $L^0=M^0=B$, l'isomorphisme précédent
est un isomorphisme de $\gr_{n+1-c}(B)$-algèbres, autrement
dit les algèbres $H^{0}(L^\bullet)$ et $H^{0}(M^\bullet)$ sont
$(n+1-c)$-isomorphes.  \end{itemize} \end{proposition}

\begin{proof} Les deux premiers points sont prouvés dans le lemme
3.1 de \textit{loc. cit.}. Le dernier est trivial. Reste le point
(iii).

Pour $n=c$, c'est le théorème 3.2  de \textit{loc. cit.}
dont on ne fait qu'adapter la preuve dans le
cas $n>c$. Soit $m\in\ZZ$. On écrit $\d_L,\d_M$
pour $\d_L^{-1},\d^{-1}_M$. Pour $\delta=\d_L,\d_M$, on a
$$\gr_{n+1-c}^m(\coker(\delta))={I}^mL^0/({I}^{m+n+1-c}L^0+{I}^mL^0\cap\Im(\delta))$$
de sorte qu'il s'agit de montrer l'égalité
$${I}^{m+n+1-c}L^0+{I}^mL^0\cap\Im(\d_L)={I}^{m+n+1-c}L^0+{I}^mL^0\cap\Im(\d_M)$$
pour tout $m\in\ZZ$. Soit $x\in L^{-1}$ tel que $\d_L(x)\in{I}^mL^0$.

Supposons $m\leq c$.  Comme $$\d_L(x)-\d_M(x)\in{I}^{n+1}L^0\text{
et }m\leq c\leq n,$$ on a $\d_L(x)-\d_M(x)\in{I}^{m}L^0$ de sorte
que $$\d_M(x)=\d_L(x)+\d_M(x)-\d_L(x)\in{I}^mL^0\cap\Im(d_M).$$
Comme $n+1\geq m+n+1-c$, on a également
$$\d_L(x)-\d_M(x)\in{I}^{n+1}L^0\subset{I}^{m+n+1-c}L^0$$
de sorte que $$\d_L(x)=\d_L(x)-\d_M(x)+\d_M(x)\in
{I}^{m+n+1-c}L^0+{I}^mL^0\cap\Im(\d_M)$$ et donc
$${I}^{m+n+1-c}L^0+{I}^mL^0\cap\Im(\d_L)\subset{I}^{m+n+1-c}L^0+{I}^mL^0\cap\Im(\d_M).$$
Par symétrie des rôles de $\d_L$ et $\d_M$, on a l'égalité
cherchée dans ce cas.

Si $m>c$, le calcul est analogue. On a (\ref{3-AR})
$${I}^mL^0\cap\Im(\d_L)\subset{I}^{m-c}\d_L(L^{-1})$$ de sorte
que $$\d_L(x)=\d_L(x')\text{ avec }x'\in{I}^{m-c}L^{-1}.$$
Comme $\d_L-\d_M=0\mod{I}^{n+1}$, la matrice de $\d_L-\d_M$
est à coefficients dans ${I}^{n+1}$ de sorte que
$$\d_L-\d_M\in{I}^{n+1}\Hom_B(L^{-2},L^{-1}).$$ On a donc
$$\d_L(x')-\d_M(x')\in{I}^{n+1}{I}^{m-c}L^{0}={I}^{n+1+m-c}L^{0}.$$
Comme
$$\d_M(x')=\d_L(x')+\d_M(x')-\d_L(x')=\d_L(x)+\d_M(x')-\d_L(x'),$$
on a d'une part $$\d_M(x')\in({I}^m
L^0+{I}^{n+1+m-c}L_0)\cap\Im(d_M)\stackrel{n\geq
c}{\subset}{I}^m L^0\cap\Im(d_M),$$ et, d'autre part,
$$\d_L(x)=\d_M(x')-\d_L(x')+\d_L(x')+\in{I}^{m+n+1-c}L^0+{I}^mL^0\cap\Im(\d_M).$$
On conclut comme plus haut  par symétrie.\end{proof}



\section{Modèles et $a$-isomorphismes}
\begin{theoreme}[Approximation]\label{3-theo} Soit $A$  un
anneau local noethérien, $\m$ son idéal maximal, $\hat A$
son complété. On suppose $A$ excellent et hensélien. Soit
$\pi:{\s}=\Spec(\hat A)\ra S=\Spec(A)$ le morphisme canonique. Soit
$\X$ de type fini sur $S$. On se donne de plus $\alpha_0\in E$
et  un modèle (cf. \ref{3-def-mod}) $X_{\alpha_0}$ de $\X$
sur $S_{\alpha_0}$.  Pour tout $\alpha\geq \alpha_0$ on note
$X_\alpha=X_\alpha\times_{S_{\alpha_0}}S_\alpha$ le modèle de
$\X$ sur $S_\alpha$ déduit par changement de base.  Il existe
$\alpha_1\geq\alpha_0$ et des entiers $n_0\geq c>0$ tels que
pour tout $n\geq n_0,\alpha\geq\alpha_1$ et toute section $u$ de
$t_\alpha$ qui est $(n+1)$-proche de $s_\alpha$, il existe un unique
$(n+1-c)$-isomorphisme $\X\iso_{n-c}X_u$ au dessus de l'isomorphisme
$X_{n-c}\ra (X_u)_{n-c}$ déduit de (\ref{3-Xn}).\end{theoreme}


\begin{definition}\label{3-rem} Dans les conditions précédentes,
on dit que $(X_\alpha,\alpha,u)$ (ou, si aucune conclusion n'est
à craindre, $X_u$) est une \emph{approximation} de $\X$ sur $S$
(à l'ordre $n-c$).  \end{definition}

L'assertion « Il existe $\alpha_0$, un entier  $n_0$ tels
que pour tout $n\geq n_0,\alpha\geq\alpha_0$ et toute section
$u$ de $t_\alpha$ qui est $(n+1)$-proche de $s_\alpha$, $X_u$
vérifie la propriété $P$ » pourra parfois être condensée
en « Toute approximation $X_u$ assez fine de $\X$ vérifie la
propriété $P$ ». On emploiera une terminologie analogue pour
les approximations de ${\s}$-morphismes.

\begin{proof}Deux $(n+1-c)$-isomorphismes diffèrent par un
automorphisme $$\iota:\gr_{n+1-c}(\OO_\X)\iso\gr_{n+1-c}(\OO_\X)$$
de $\OO_{X_{n-c}}$-algèbres graduées. Il est en particulier
$\OO_S$-linéaire. Comme $\gr_{n+1-c}(\OO_\X)$ est engendré sur
$\gr_a(\OO_S)$ par $\OO_{X_{n-c}}$, l'automorphisme $\iota$ est
l'identité. D'où l'unicité.

On peut donc supposer $\X$ affine. Comme $\X$ est de type fini sur
$\s$, $X$ se plonge dans l'espace affine $$\Aff^m_{\s}=\Spec(\hat
A[t])$$ de coordonnées $t=(t_1,\cdots,t_m)$ comme le sous-schéma
fermé d'idéal $$J=\langle\tilde P_1,\cdots,\tilde P_N\rangle$$ où
$\tilde P_i\in B=\hat A[t]$. Choisissons une résolution partielle
du $B$-module $C=B/J$ par des $B$-modules libres de type fini

 \begin{equation}\label{3-1}
    B^a\xrightarrow{\tilde R}B^b\xrightarrow{\tilde P=(\tilde
    P_i)}B\ra C\ra 0
\end{equation} où $\tilde R$ est une matrice à coefficients
dans $B$.

Pour $\alpha_0$ assez grand,  $\tilde P$  et $\tilde R$ proviennent
de matrices $P_{\alpha_0},R_{\alpha_0}$ à coefficients dans
$$ B_{\alpha_0}=A_{\alpha_0}[t]\text{ telles que }PR=0$$ de
sorte que le fermé $F$ de $\Aff^m_{A_{\alpha_0}}$ d'équations
$P_{\alpha_0,1}=\cdots=P_{\alpha_0,N}=0$ est un modèle de $\X$
sur $S_{\alpha_0}$. Comme rappelé dans la section~\ref{3-mod},
quitte à changer $\alpha_0$ en un indice plus grand, on peut
supposer qu'on a $F=X_{\alpha_0}$. Pour $\alpha\geq\alpha_0$, note
$P_\alpha,R_\alpha$ les matrices à coefficients dans $B_\alpha$
déduites de $P_{\alpha_0},\R_{\alpha_0}$ par le morphisme
$$B_{\alpha_0}=A_{\alpha_0}[t]\ra B_\alpha=A_\alpha[t].$$
Pour tout $\alpha\geq \alpha_0$, les matrices à coefficients
dans $B$ déduites de $P_{\alpha},\R_{\alpha}$ par le morphisme
$$B_{\alpha_0}=A_{\alpha_0}[t]\ra B=\hat A[t]$$ sont les mêmes :
on les note $P,R$.

On s'est ramené, pour $\alpha\geq \alpha_0$,
au cas où $$X_\alpha=\Spec(C_\alpha)\text{ avec
}C_\alpha=B_\alpha/(P_\alpha).$$ On dispose donc d'une part
d'un complexe (en degrés [-2,0]) de $B_\alpha$-modules libres
\begin{equation*}
  L_\alpha=(
  B_\alpha^a\xrightarrow{R_\alpha}B_\alpha^b\xrightarrow{P_\alpha=(P_{i,\alpha})}B_\alpha)
\end{equation*} avec $H^0(L_\alpha)=C_\alpha$. Le complexe de
$B$-modules libres de rang fini \begin{equation*}
  L=B\otimes_{B_\alpha}L_\alpha=(
  B^a\xrightarrow{R}B^b\xrightarrow{P=1\otimes P_\alpha}B)
\end{equation*}  est acyclique en degré $-1$ par construction.

\begin{remarque}\textit{A priori}, $L_\alpha$ n'a pas de
raison d'être acyclique en degré $-1$, même pour $\alpha$
grand. \end{remarque}


D'autre part, la section $u$ de $t_\alpha$ est définie par un
morphisme de $A$-algèbres $$u^*:A_\alpha\ra A$$ de sorte que
$$u^*\mod\m^{n+1}=s^*_\alpha\mod\hat\m^{n+1},$$
 où $s_\alpha^*: A_\alpha\ra\hat A$ est défini par $s_\alpha:\hat
 S\ra S_\alpha$ (\ref{3-salpha}). Par action sur les coefficients
 des polynômes, on obtient un morphisme d'anneau
$$\bar u:B_\alpha=A_\alpha[t]\ra A[t]\ra\hat A[t]=B$$ d'où un
complexe \begin{equation*}
  M=(  B^a\xrightarrow{\bar u(R)}B^b\xrightarrow{\bar u(P)}B)
\end{equation*}  Par construction, on a $$L/\m^{n+1}L=M/\m^{n+1}M.$$
On choisit alors une constante d'Artin-Rees $c$ pour
$B^{b}\xrightarrow{P}B$ et on invoque la proposition \ref{3-CdJ}
pour conclure.  \end{proof}



\begin{corollaire}\label{3-perm}Soient $\X,{{\Y}}$ des
$S$-schémas noethériens qui sont $a$-proches. Soit $x\in X_0=Y_0$.
\begin{itemize} \item[(i)] Si $a\geq 1$, les dimensions de $\X$ et
${{\Y}}$ en $x$ sont les mêmes.  \item[(ii)] Si $a\geq 2$ et $\X$
régulier en $x$, alors ${{\Y}}$ régulier en $x$.  \item[(iii)]
Supposons $\X\ra{\s}$ de type fini et $\X$ régulier. Alors,
il existe $\alpha_0\in E, n_0\in\mathbf{N}$ tels que  pour tout
$\alpha\geq\alpha_0$, tout entier $n\geq n_0$, toute section $u$
de $t_\alpha$ qui est $n$-proche de $s_\alpha:{\s}\ra S_\alpha$,
tout modèle $X_\alpha$ de $\X$ sur $S_\alpha$, le schéma $X_u$
soit régulier dans un voisinage ouvert de la fibre spéciale.
\end{itemize} \end{corollaire}


\begin{proof}

Par hypothèses, les cônes normaux de $X_0,Y_0$ dans $\X,{{\Y}}$
sont $S$-isomorphes. Comme la dimension de $\X$ en $x$ est égale
à celle de son cône normal \cite[15.9]{CRT@Matsumura}, le premier
point en découle.

Supposons maintenant que $\X,{{\Y}}$ soient $2$-proches. D'après
(i), on sait que $\X$ et ${{\Y}}$ ont même dimension en $x$. Comme
$\X,{{\Y}}$ sont $2$-proches, $X_1$ et $Y_1$ sont isomorphes. Puisque
l'espace tangent de Zariski à $\X$ en un point de $X_0$ ne dépend
que de son second voisinage infinitésimal $X_1$,  les $k(x)$-espaces
vectoriels cotangents de Zariski en $x$ à $\X$ et ${{\Y}}$ sont
isomorphes, d'où ii).

Pour le dernier point, il suffit d'invoquer les deux premiers et
le théorème~\ref{3-theo} pour conclure qu'une approximation assez
fine est régulière au voisinage de la fibre spéciale. Comme $X_u$
est excellent (puisque de type fini sur $S$ excellent), son lieu
régulier $R$ est ouvert de sorte que $R$ est un voisinage ouvert
régulier de la fibre spéciale.\end{proof}

 \begin{remarque}\label{3-remgabber} O. Gabber sait généraliser
 la proposition~\ref{3-CdJ} au cas où les complexes envisagés
 sont seulement de type fini sur un anneau noethérien pour
 obtenir  les proximités de la cohomologie également en
 degré $-2$ (et pas seulement en degré $0,-1$). Il peut plus
 précisément montrer des énoncés de proximité pour les images,
 noyaux des différentielles\footnotemark. Gabber en déduit de
 nombreux énoncés de permanence par approximation analogues au
 corollaire~\ref{3-perm} . Notamment, si $X,Y$ sont $a$-proches pour
 $a$ assez grand, alors $X$ réduit (resp. normal) le long de $X_0$ entraîne $Y$
 réduit (resp. normal) le long de $Y_0$. Cependant, plusieurs questions naturelles
 restent suspens comme par exemple la permanence des propriétés
 $S_n,R_n$. \end{remarque}
\footnotetext{La preuve de cette généralisation a été exposée
par A. Moreau lors du séminaire sur les travaux de Gabber.}
\section{Réduction au cas local noethérien complet}Rappelons
l'énoncé du théorème d'uniformisation

\begin{theoreme}[Uniformisation]\label{3-unif} Soient $T$ un schéma
noethérien quasi-excellent et $Z$ un fermé rare de $T$. Soit $\ell$
un nombre premier inversible sur $T$.  Il existe une famille finie de
morphismes $(X_i\ra T)_{i\in I}$  et telle que pour tout $i\in I$ on
ait \begin{itemize} \item[(i)]La famille finie de morphismes $(X_i\ra
T)_{i\in I}$   est {\pspf}-couvrante (resp. \altl-couvrante) ;
  \item[(ii)] $X_i$ est régulier et connexe ; \item[(iii)]l'image
  inverse de $Z$ dans $X_i$ est vide ou le support d'un diviseur
  à croisements normaux stricts;
\end{itemize} \end{theoreme}



Nous allons montrer l'énoncé de réduction suivant.

 \begin{proposition}\label{3-reduc} Si (\ref{3-unif}) est vrai pour
 tout $T$ noethérien, local, complet, alors (\ref{3-unif}) est vrai.
  \end{proposition}
\begin{proof} On peut d'abord supposer  $T$ local, excellent et
hensélien (rappelons (\ref{1-cor-normal-versus-complet}) qu'un
schéma local, hensélien et quasi-excellent est excellent).

En effet, supposons le théorème prouvé dans ce cas. D'après
\ega{IV$_4$}{18.7.6}, l'hensélisé $T_{(t)}$ de $T$ en $t\in T$
est excellent et $Z_{(t)}$ est un fermé rare de $T_{(t)}$. On peut
alors trouver une famille finie $(f_i:X_i\ra T_{(t)})_{i\in I}$
vérifiant les 3 propriétés (i)-(iii) du théorème. La famille
$f_i$ est dominée par un recouvrement standard $f:Y\ra T'\ra
T_{(t)}$ avec $Y\ra T'$ Zariski (resp. Nisnevich) couvrant et
$T'\ra T$ propre et surjectif (resp. propre et surjectif tels que
pour tout point maximal $\eta$ de $T$, il existe un point maximal
$\eta'$ de $T'$ au dessus de $\eta$ avec  $\ell$ ne divisant pas
$\deg(k(\eta')/k(\eta))$). D'après \ega{IV$_3$}{8.2.2}, les $f_i$
et $f$ proviennent de morphismes $\tilde f_i:\tilde X_i\ra \tilde
T$ où  $\tilde T\ra T$ est un voisinage Nisnevich étale de $T$ en
$t$. La topologie de Nisnevich étant moins fine que les topologies
$\pspf$ et $\altl$ (\refext{II}{2-altération plus fine que étale
cd} (ii)), les propriétés usuelles de passage à la limite
(\ref{3-stabmod} et\ega{IV$_3$}{8.4.3}, \ega{IV$_4$}{19.8.1})
permettent de conclure comme plus haut.

Supposons donc $T$, schéma local, hensélien, excellent.

Quitte à remplacer $S$ par la somme disjointe de ses composantes
réduites, on se ramène au cas où $S$ est de plus intègre.

On peut supposer de plus $S=\Spec(A)$ normal intègre. En
effet, comme $A$ est excellent, le morphisme de normalisation
est fini de degré générique $1$, donc est $\pspf$-couvrant
(resp. $\altl$-couvrant). Comme $A$ est local intègre et hensélien,
$A$ est unibranche de sorte que le normalisé de $A$ est local,
donc intègre puisque normal et est noethérien hensélien puisque
fini sur $A$.

Comme $A$ est excellent, la normalisation commute à la complétion
(\ref{1prop-normal-versus-complet}) de sorte que $\hat A$ est dès
lors normal comme $A$, donc également intègre puisque normal
et local.

On peut donc supposer $T=S$ avec $S$ schéma local intègre, normal,
hensélien et excellent.

Comme $\s$ est plat sur $S$, l'image inverse ${\Z}$ de $Z$ est
encore un fermé rare de ${\s}$.  Choisissons une uniformisation
$$(\tilde\X_i\ra {\s})_{i\in I'}$$ de $({\s},{\Z})$ comme
dans ~\ref{3-unif}.  D'après ~\ref{3-theo},~\ref{3-alt-mod}
et \ref{3-perm}, on peut trouver un entier $n\geq 0$, des
$n$-isomorphismes $\tilde\X_i\ra_{n}(\tilde X_i)_u$  de sorte  que

 a)  chaque $S$-schéma $(X_i)_u$ est régulier le long de sa fibre
 spéciale $(\tilde X_i)_0$, donc au voisinage (le lieu régulier
 étant ouvert puisque les schémas considérés sont excellents).

 b) la famille $((\tilde X_i)_u)_\r$ est \pspf-couvrante
 (resp. \altl-couvrante).

 D'après a), $(\tilde X_i)_u$ est régulier au voisinage de la
 fibre spéciale et y est la réunion disjointe de ses composantes
 connexes qui sont intègres. Ainsi, au voisinage de la fibre
 spéciale, $((\tilde X_i)_u)_\r$ est schématiquement la réunion
 disjointe des composantes de $(\tilde X_i)_u$ dominant $S$. Comme
 tout voisinage ouvert de la fibre spéciale $(\tilde X_i)_0$ dans
 $((\tilde X_i)_u)_\r$ est {\pspf}-couvrant (resp. \altl-couvrant)
 (\refext{II}{2-voisinages fibre restent couvrants}), la famille
 $(X_i\ra S)_{i\in I}$ des composantes connexes de voisinages
 convenables des $(\tilde X_i)_0$ dans  $(\tilde X_i)_u, i\in I'$
 vérifie les conditions (i) et (ii).


 Soit $D'$ l'image inverse de ${\Z}$ dans $\X=\sqcup_{i\in I}
 X_i$ qu'on peut supposer non vide. Par hypothèse, $D=D'_\red$
 est un diviseur à croisements normaux strict, c'est-à-dire
 $D=\sum_{j\in J} D_j$ avec  $$D_K=\cap_{j\in K} D_j$$ régulier de
 codimension $\text{card}(J)$ pour toute partie $K\subset J$. Quitte
 à augmenter $\alpha$, on peut supposer que les $D_j$ ont des
 modèles sur $S_\alpha$, ces modèles induisant des modèles des
 $D_J$. Comme $u$ est une section de $t_\alpha$, le schéma $D_u$
 réunion schématique des $(D_i)_u$ est, topologiquement, l'image
 inverse de $Z$ dans $X_u$. D'après ~\ref{3-perm}, on peut supposer
 que chaque $(D_u)_K$ est régulier de codimension $\text{card}(K)$
 (par caténarité de $X_i$ qui est excellent puisque de type fini
 sur $S$) de sorte que $D_u$ est un diviseur à croisements normaux
 stricts le long de la fibre spéciale. Les lieux réguliers de
 $(D_u)_K$ et $X_u$ étant ouverts, on peut supposer que $D_u$
 est un diviseur à croisements normaux stricts au voisinage de la
 fibre spéciale (caténarité de $X_u$).
\end{proof}


\ifx\danslelivre\undefined
\bibliography{../fichiers-configuration-latex/bibliographie}
\bibliographystyle{../fichiers-configuration-latex/smfalpha}


\end{document} \fi
